% GENERATED FILE. Edit canonical lessons, then run npm run generate:books.
% canonical-source-hash: fnv1a64:6598048916680ca2
% canonical-lessons: TA-C386-utti, TA-C386-munneccarikkai, TA-C386-varaivu, TA-C386-munnetu, TA-C386-kurikkol

\chapter{Planning Ahead}
\label{ch:ta-planning-ahead}
\hlchaptermodality{386}

\begin{chapteropening}
\textbf{By the end of this chapter:} \emph{I can talk about a strategy, a precaution and a goal.}

Complete the last lesson of chapter 386: I can talk about a strategy, a precaution and a goal.
\end{chapteropening}

\section[\texorpdfstring{\ta{உத்தி} (utti)}{உத்தி (utti)}]{\ta{உத்தி} (utti) — a strategy; a technique}

\label{lesson:TA-C386-utti}



\begin{quote}
Before the new one: say the Tamil for a complication, then the Tamil for to handle.
\end{quote}

\subsection*{You'll want to know: \ta{உத்தி}}
\textbf{\ta{உத்தி}} — \emph{utti} — \textquotedblleft{}a strategy; a technique\textquotedblright{}.

\begin{quote}
\ta{வெற்றி} \ta{பெற} \ta{நமக்கு} \ta{ஒரு} \ta{புதிய} \ta{உத்தி} \ta{தேவை}.
\end{quote}

\begin{quote}
\emph{veṟṟi peṟa namakku oru putiya utti tēvai.}
\end{quote}

\begin{quote}
We need a new strategy to win.
\end{quote}

Often plural: \ta{உத்திகள்}, techniques.

\begin{cousinweb}
\textbf{\ta{உத்தி}} (\emph{utti}) is borrowed from Sanskrit \emph{yukti}, \textquotedblleft{}union, preparation, application\textquotedblright{} (CDIAL 10480); Tamil drops the initial y and makes kt into \ta{த்த}.
\end{cousinweb}

\begin{grammarlens}[title={What you've built}]
One more word for saying what happens if.
\end{grammarlens}

\subsection*{Your turn}
\begin{itemize}
\raggedright
  \item \emph{Say it:} \emph{utti}
  \item \emph{Say it:} \emph{utti}, once more
  \item \emph{Say it:} \emph{veṟṟi peṟa namakku oru putiya utti tēvai.} — We need a new strategy to win.
  \item \emph{Recall:} say the Tamil for a complication, then the Tamil for to handle, then say \emph{utti} again
\end{itemize}

\begin{tcolorbox}[breakable,colback=teal!4,colframe=teal!35!black,title={Before you move on}]
What does \ta{உத்தி} mean? (\textquotedblleft{}A strategy; a technique\textquotedblright{}.) Say it once more.
\end{tcolorbox}
\section[\texorpdfstring{\ta{முன்னெச்சரிக்கை} (muṉṉeccarikkai)}{முன்னெச்சரிக்கை (muṉṉeccarikkai)}]{\ta{முன்னெச்சரிக்கை} (muṉṉeccarikkai) — a precaution}

\label{lesson:TA-C386-munneccarikkai}



\begin{quote}
Before the new one: say the Tamil for to handle, then the Tamil for a strategy.
\end{quote}

\subsection*{You'll want to know: \ta{முன்னெச்சரிக்கை}}
\textbf{\ta{முன்னெச்சரிக்கை}} — \emph{muṉṉeccarikkai} — \textquotedblleft{}a precaution\textquotedblright{}.

\begin{quote}
\ta{மழைக்காலத்துக்கு} \ta{முன்} \ta{முன்னெச்சரிக்கை} \ta{எடுக்க} \ta{வேண்டும்}.
\end{quote}

\begin{quote}
\emph{maḻaikkālattukku muṉ muṉṉeccarikkai eṭukka vēṇṭum.}
\end{quote}

\begin{quote}
Precautions must be taken before the rainy season.
\end{quote}

\ta{முன்}, before, + a warning: \ta{எச்சரி}, to warn, is the verb inside it.

\begin{cousinweb}
\textbf{\ta{முன்னெச்சரிக்கை}} (\emph{muṉṉeccarikkai}) is built from native \emph{muṉ}, \textquotedblleft{}before\textquotedblright{} (DEDR 5020), and \emph{eccarikkai}, \textquotedblleft{}a warning\textquotedblright{}, the noun of \ta{எச்சரி}: a warning taken ahead of time.
\end{cousinweb}

\begin{grammarlens}[title={What you've built}]
One more word for saying what happens if.
\end{grammarlens}

\subsection*{Your turn}
\begin{itemize}
\raggedright
  \item \emph{Say it:} \emph{muṉṉeccarikkai}
  \item \emph{Say it:} \emph{muṉṉeccarikkai}, once more
  \item \emph{Say it:} \emph{maḻaikkālattukku muṉ muṉṉeccarikkai eṭukka vēṇṭum.} — Precautions must be taken before the rainy season.
  \item \emph{Recall:} say the Tamil for to handle, then the Tamil for a strategy, then say \emph{muṉṉeccarikkai} again
\end{itemize}

\begin{tcolorbox}[breakable,colback=teal!4,colframe=teal!35!black,title={Before you move on}]
What does \ta{முன்னெச்சரிக்கை} mean? (\textquotedblleft{}A precaution\textquotedblright{}.) Say it once more.
\end{tcolorbox}
\section[\texorpdfstring{\ta{வரைவு} (varaivu)}{வரைவு (varaivu)}]{\ta{வரைவு} (varaivu) — a draft}

\label{lesson:TA-C386-varaivu}



\begin{quote}
Before the new one: say the Tamil for a strategy, then the Tamil for a precaution.
\end{quote}

\subsection*{You'll want to know: \ta{வரைவு}}
\textbf{\ta{வரைவு}} — \emph{varaivu} — \textquotedblleft{}a draft\textquotedblright{}.

\begin{quote}
\ta{அவர்} \ta{கடிதத்தின்} \ta{வரைவைத்} \ta{தயாரித்தார்}.
\end{quote}

\begin{quote}
\emph{avar kaṭitattiṉ varaivait tayārittār.}
\end{quote}

\begin{quote}
He prepared a draft of the letter.
\end{quote}

\ta{வரைவுத்} \ta{திட்டம்}, a draft plan.

\begin{cousinweb}
\textbf{\ta{வரைவு}} (\emph{varaivu}) is native Dravidian: DEDR 5263 lists Tamil \emph{varai}, \textquotedblleft{}to write, inscribe, paint, draw\textquotedblright{}, and \emph{varaivu}: a first drawing of a text.
\end{cousinweb}

\begin{grammarlens}[title={What you've built}]
One more word for saying what happens if.
\end{grammarlens}

\subsection*{Your turn}
\begin{itemize}
\raggedright
  \item \emph{Say it:} \emph{varaivu}
  \item \emph{Say it:} \emph{varaivu}, once more
  \item \emph{Say it:} \emph{avar kaṭitattiṉ varaivait tayārittār.} — He prepared a draft of the letter.
  \item \emph{Recall:} say the Tamil for a strategy, then the Tamil for a precaution, then say \emph{varaivu} again
\end{itemize}

\begin{tcolorbox}[breakable,colback=teal!4,colframe=teal!35!black,title={Before you move on}]
What does \ta{வரைவு} mean? (\textquotedblleft{}A draft\textquotedblright{}.) Say it once more.
\end{tcolorbox}
\section[\texorpdfstring{\ta{முன்னெடு} (muṉṉeṭu)}{முன்னெடு (muṉṉeṭu)}]{\ta{முன்னெடு} (muṉṉeṭu) — to take forward; to spearhead}

\label{lesson:TA-C386-munnetu}



\begin{quote}
Before the new one: say the Tamil for a precaution, then the Tamil for a draft.
\end{quote}

\subsection*{You'll want to know: \ta{முன்னெடு}}
\textbf{\ta{முன்னெடு}} — \emph{muṉṉeṭu} — \textquotedblleft{}to take forward; to spearhead\textquotedblright{}.

\begin{quote}
\ta{இந்த} \ta{முயற்சியை} \ta{அவர்தான்} \ta{முன்னெடுத்தார்}.
\end{quote}

\begin{quote}
\emph{inta muyaṟciyai avartāṉ muṉṉeṭuttār.}
\end{quote}

\begin{quote}
It was he who took this effort forward.
\end{quote}

Past (I …): \ta{முன்னெடுத்தேன்}, \emph{muṉṉeṭuttēṉ}.

What is led takes -\ta{ஐ}: \ta{போராட்டத்தை} \ta{முன்னெடு}, to lead a struggle.

\begin{cousinweb}
\textbf{\ta{முன்னெடு}} (\emph{muṉṉeṭu}) is native: \emph{muṉ}, \textquotedblleft{}front\textquotedblright{} (DEDR 5020), and \emph{eṭu}, \textquotedblleft{}to take up, lift\textquotedblright{} (DEDR 851): to take something to the front.
\end{cousinweb}

\begin{grammarlens}[title={What you've built}]
One more word for saying what happens if.
\end{grammarlens}

\subsection*{Your turn}
\begin{itemize}
\raggedright
  \item \emph{Say it:} \emph{muṉṉeṭu}
  \item \emph{Say it:} \emph{muṉṉeṭu}, once more
  \item \emph{Say it:} \emph{inta muyaṟciyai avartāṉ muṉṉeṭuttār.} — It was he who took this effort forward.
  \item \emph{Recall:} say the Tamil for a precaution, then the Tamil for a draft, then say \emph{muṉṉeṭu} again
\end{itemize}

\begin{tcolorbox}[breakable,colback=teal!4,colframe=teal!35!black,title={Before you move on}]
What does \ta{முன்னெடு} mean? (\textquotedblleft{}To take forward; to spearhead\textquotedblright{}.) Say it once more.
\end{tcolorbox}
\section[\texorpdfstring{\ta{குறிக்கோள்} (kuṟikkōḷ)}{குறிக்கோள் (kuṟikkōḷ)}]{\ta{குறிக்கோள்} (kuṟikkōḷ) — an aim; a goal}

\label{lesson:TA-C386-kurikkol}



\begin{quote}
Before the new one: say the Tamil for a draft, then the Tamil for to take forward.
\end{quote}

\subsection*{You'll want to know: \ta{குறிக்கோள்}}
\textbf{\ta{குறிக்கோள்}} — \emph{kuṟikkōḷ} — \textquotedblleft{}an aim; a goal\textquotedblright{}.

\begin{quote}
\ta{என்} \ta{குறிக்கோள்} \ta{ஒரு} \ta{மருத்துவர்} \ta{ஆவது}.
\end{quote}

\begin{quote}
\emph{eṉ kuṟikkōḷ oru maruttuvar āvatu.}
\end{quote}

\begin{quote}
My aim is to become a doctor.
\end{quote}

\ta{குறி}, a mark, a target + \ta{கோள்}, a holding.

\begin{cousinweb}
\textbf{\ta{குறிக்கோள்}} (\emph{kuṟikkōḷ}) is native: \emph{kuṟi}, \textquotedblleft{}a mark; to aim at\textquotedblright{} (DEDR 1847), and \emph{kōḷ}, \textquotedblleft{}taking, holding\textquotedblright{} (DEDR 2151): the mark one holds to.
\end{cousinweb}

\begin{grammarlens}[title={What you've built}]
One more word for saying what happens if.
\end{grammarlens}

\subsection*{Your turn}
\begin{itemize}
\raggedright
  \item \emph{Say it:} \emph{kuṟikkōḷ}
  \item \emph{Say it:} \emph{kuṟikkōḷ}, once more
  \item \emph{Say it:} \emph{eṉ kuṟikkōḷ oru maruttuvar āvatu.} — My aim is to become a doctor.
  \item \emph{Recall:} say the Tamil for a draft, then the Tamil for to take forward, then say \emph{kuṟikkōḷ} again
\end{itemize}

\begin{tcolorbox}[breakable,colback=teal!4,colframe=teal!35!black,title={Before you move on}]
What does \ta{குறிக்கோள்} mean? (\textquotedblleft{}An aim; a goal\textquotedblright{}.) Say it once more.
\end{tcolorbox}
